\documentclass[11pt]{article}
\usepackage[margin=1in]{geometry}
\usepackage{amsmath,amssymb,amsthm}
\usepackage[colorlinks=true,urlcolor=blue,citecolor=blue,linkcolor=blue]{hyperref}
\newtheorem{theorem}{Theorem}
\newtheorem{lemma}[theorem]{Lemma}
\newtheorem{corollary}[theorem]{Corollary}
\newcommand{\R}{\mathbb R}
\newcommand{\norm}[1]{\left\lVert#1\right\rVert}
\newcommand{\MM}{T_{\rm MM}}
\newcommand{\tO}{\widetilde O}
\newcommand{\rank}{\operatorname{rank}}
\newcommand{\range}{\operatorname{range}}
\newcommand{\diag}{\operatorname{diag}}
\title{Sketch-preconditioned thresholds\\for spectral approximation of a factored product}
\author{Research candidate: independent novelty audit in progress}
\date{November 14, 2026}
\begin{document}
\maketitle
\begin{abstract}
We develop algorithms that trade factor-sketch accuracy against the cost
of a preconditioned spectral threshold test. A two-level variant retains
the finer sketches as operators, avoiding their dense Gram matrices.
For \(n=p=N\), \(m=\lceil N^{50/59}\rceil\), and
\(\theta=m^{-56/125}\), a fixed-depth extension constructs a rank-one approximation in
\(o(N^2)\) arithmetic operations and explicitly expands it in total
\(O(N^2)\) time. The error is at most \((1+\theta)\sigma_2(AB)\)
with any fixed prescribed confidence, for every fixed real factor pair.
Schur-complement deflation extends the guarantee to general rank, and a
stopping rule certifies absolute error without a supplied target rank.
The strongest strict quadratic-time consequence holds for fixed rank.
The bounds improve the audited general-product Gram and Krylov bounds;
global novelty remains under review. The original \(\theta=1/m\) target
and finite-precision stability remain open.
\end{abstract}

\section{Statement and arithmetic model}
Let \(A\in\R^{n\times m}\), \(B\in\R^{m\times p}\),
\(m\leq\min(n,p)\), \(M=AB\), and \(I=m(n+p)\).
The algorithm returns \(u\in\R^n\), \(v\in\R^p\), and explicitly
forms \(C=uv^\top\). All costs below are real arithmetic operations;
square roots, inner products, comparisons, and standard numerical small-core
eigenvalue computations are included in the usual arithmetic linear-algebra
model. A stable floating-point/bit-complexity theorem is not claimed.
Exponent slack \(\zeta>0\) absorbs fast-algebra overhead. Logarithms in
dimensions, inverse accuracy, and inverse failure probability are suppressed.
Singular values and eigenvalues beyond the matrix dimension are defined as zero.
All query counts are capped in advance. Test required squared norms and
denominators before normalization or inversion; a nonpositive required norm,
singular auxiliary solve, or invalid bracket returns the current factors.
These exits are absent on the good events proved below. They keep the
algorithm well defined with bounded runtime on failure outcomes as well.

For \(0<h\leq1/16\), let \(s_A,s_B\) be the dimensions of fast
squared-norm \((1\pm h)\) embeddings for \(\range A\) and
\(\range B^\top\), with joint failure probability at most \(\delta/4\).
SRHT embeddings permit
\[
 s_A=\min\{n,\tO(m/h^2)\},\qquad
 s_B=\min\{p,\tO(m/h^2)\},
\]
using the identity in a smaller ambient dimension. Padding to a power of two
changes only constants. Define
\[
 G_h=\MM(m,s_A,m)+\MM(m,s_B,m).
\]

\begin{theorem}[Candidate rank-one bound]\label{thm:main}
For \(0<\theta\leq h\leq1/16\) and \(0<\delta<1\), the
algorithm below satisfies, for every fixed \(A,B\),
\[
 \Pr\{\norm{M-C}_2\leq(1+\theta)\sigma_2(M)\}\geq1-\delta
\]
with total arithmetic cost
\begin{equation}\label{eq:cost}
 \tO\!\left(I+G_h+m^{\omega+\zeta}
              +\frac{m^2}{\sqrt h}
              +I\sqrt{h/\theta}+np\right).
\end{equation}
The formula includes rank-deficient inputs and explicit output expansion.
\end{theorem}

The two-level bound in Theorem~\ref{thm:twolevel} strengthens this tradeoff.
Corollaries~\ref{cor:twolevel} and~\ref{cor:multilevel} give increasingly
accurate quadratic-output regimes.

The proof uses standard randomized Krylov eigenvalue estimation \cite{jyz}
and fast subspace embeddings \cite{cnw}. Preconditioned Lanczos for shifted
generalized eigenvalue pencils is classical \cite{templates}; it is not claimed
as a new mechanism. The potentially new contribution is the two-factor
sketch construction, the global error-to-threshold accounting, and its optimized
total cost. Sketch-preconditioned eigenvector refinement also has recent prior
art, including EPSI \cite{epsi}; novelty of this bound is not yet settled.

\section{Implicit coordinates, including rank deficiency}
Write \(P=A^\top A\), \(Q=BB^\top\). Construct
\(\widetilde P=(SA)^\top SA\) and
\(\widetilde Q=(TB^\top)^\top TB^\top\), so on the good event
\begin{equation}\label{eq:ose}
 (1-h)P\preceq\widetilde P\preceq(1+h)P,\qquad
 (1-h)Q\preceq\widetilde Q\preceq(1+h)Q.
\end{equation}
Let \(a=\rank\widetilde P=\rank P\). If \(a=0\), return zero.
Small-core rank-revealing algebra supplies a full-column-rank matrix
\(W\in\R^{m\times a}\) such that
\(\range W=\range P\) and \(W^\top\widetilde P W=I_a\).
Set \(L=W^\dagger=(W^\top W)^{-1}W^\top\). No exact spectral
oracle is needed for this support reduction: choose independent columns
\(W_0\) of \(\widetilde P\) by rank-revealing elimination, form the
positive definite core \(W_0^\top\widetilde P W_0=R^\top R\), and set
\(W=W_0R^{-1}\). These products, inversions, and Cholesky operations
cost \(\tO(m^{\omega+\zeta})\) in the stated arithmetic model.

The following objects are used implicitly:
\[
 Z=AW,\quad G=Z^\top Z=W^\top P W,\quad
 Q'=LQL^\top,\quad Q_0=L\widetilde QL^\top.
\]
Since \(WL\) is the orthogonal projector onto \(\range P\),
\(A=AWL\), so \(M=ZLB\) exactly. In particular,
\[
 \frac1{1+h}I_a\preceq G\preceq\frac1{1-h}I_a,
 \qquad (1-h)Q'\preceq Q_0\preceq(1+h)Q'.
\]
The spectra of \(H=MM^\top\) and \(G^{1/2}Q'G^{1/2}\) agree
after padding the latter with \(n-a\) additional zeros.
Set \(\lambda_1=\norm H_2\).
If \(Q_0=0\), then \(M=0\) on the good embedding event, so return zero.
Every application of \(G\) or \(Q'\) takes \(O(I+m^2)=O(I)\)
operations through the original factors. The algorithm never forms \(Z\)
or \(LB\) with all \(a\) columns/rows.

\section{A uniformly bounded threshold problem}
Let \(\nu_0=\lambda_{\max}(Q_0)>0\). A standard randomized core
Krylov computation obtains a Rayleigh value \(\nu\) satisfying
\((1-h)\nu_0\leq\nu\leq\nu_0\), with failure probability at most
\(\delta/4\), at cost \(\tO(m^2/\sqrt h)\). If the Krylov basis
saturates, use its small dense eigenproblem.
The unknown scale \(\nu_0\) is not required by the algorithm: run Lanczos
on \(Q_0\), whose spaces and Ritz vectors agree with those of
\(-Q_0\) after the sign change. Lemma 1 of \cite{jyz}, analyzed with
norm bound \(\nu_0\) and additive error \(h\nu_0\), gives an iteration
count depending only on \(h\), dimension and confidence. The scale cancels.
The comparison \eqref{eq:ose} gives
\((1-h)^2\lambda_1\leq\nu_0\leq(1+h)^2\lambda_1\).
Consequently
\begin{equation}\label{eq:bracket}
 \ell_0=\frac{\nu}{(1+h)^2}\leq\lambda_1\leq
 b_0=\frac{\nu}{(1-h)^3}.
\end{equation}
For all \(t\in[\ell_0,b_0]\), direct scalar estimates, using
\(h\leq1/16\), yield
\[
 \tfrac23\lambda_1\leq t\leq\tfrac32\lambda_1,
 \quad \norm{Q'}_2\leq2t,
 \quad \nu_0/t\leq (1+h)^2/(1-h)\leq1+4h.
\]

For such a threshold \(t\), define
\begin{equation}\label{eq:stiffness}
 K_t=tG-GQ'G,\qquad
 F_t=(1+32h)tI_a-Q_0,\qquad
 \Delta_t=F_t-K_t.
\end{equation}
Here \(K_t\) can be indefinite. The preconditioner \(F_t\) is explicit
and positive definite for every threshold in the initial bracket.

\begin{lemma}[Uniform preconditioner bounds]\label{lem:precond}
For every \(t\in[\ell_0,b_0]\),
\[
 \norm{K_t-(tI_a-Q_0)}_2\leq10ht,\qquad
 22htI_a\preceq\Delta_t\preceq42htI_a,\qquad
 F_t\succeq28htI_a.
\]
If \(F_t=J_tJ_t^\top\) is a Cholesky factorization, then
\(E_t=J_t^{-1}K_tJ_t^{-\top}\) has spectrum in \([-1,1]\).
\end{lemma}
\begin{proof}
Put \(e=h/(1-h)\leq16h/15\). The preceding bounds give
\begin{align*}
 \norm{K_t-(tI_a-Q_0)}_2
 &\leq t\norm{G-I_a}_2+
          \norm{GQ'G-Q'}_2+\norm{Q'-Q_0}_2\\
 &\leq te+(2e+e^2+h)\norm{Q'}_2\leq10ht.
\end{align*}
The claimed bounds on \(\Delta_t\) follow by adding \(32htI_a\).
The bound on \(F_t\) follows from \(\nu_0\leq(1+4h)t\).
Finally,
\[
 E_t=I_a-J_t^{-1}\Delta_tJ_t^{-\top},\qquad
 0\preceq J_t^{-1}\Delta_tJ_t^{-\top}\preceq\tfrac{42}{28}I_a.
\]
Thus its spectrum is even contained in \([-1/2,1]\).
\end{proof}

The preconditioner is allowed to depend on the threshold. A factorization
and inverse of its triangular Cholesky factor cost
\(\tO(m^{\omega+\zeta})\) per threshold, with only logarithmically
many thresholds below. A matvec with \(E_t\) takes \(O(I)\)
operations, using the implicit formula for \(K_t\) and small-core inverse
applications. No square root of the true mass matrix \(G\) is needed.
If an unexpected Cholesky failure or zero normalization occurs, the algorithm
returns zero; neither occurs on the good events used in the proof. The number
of iterations is bounded in advance, so failures do not create an unbounded
runtime outside those events.

\section{Global threshold testing and a Rayleigh vector}
We use the standard randomized Krylov fact that a fixed symmetric matrix
with spectrum in \([-1,1]\) admits, with failure probability \(\beta\),
a unit vector \(y\) whose \emph{actual} Rayleigh value satisfies
\begin{equation}\label{eq:ritz}
 \mu_{\min}(E_t)\leq q=y^\top E_ty\leq\mu_{\min}(E_t)+\tau
\end{equation}
using \(\tO(\tau^{-1/2})\) matvecs. One can apply ordinary PSD
Krylov iteration to \(I_a-E_t\succeq0\), whose norm is at most 2.
Lemma 1 of \cite{jyz} gives this actual-witness additive guarantee directly;
equivalent Chebyshev and per-vector analyses are in \cite{mm}.
The dense Ritz problem need only be solved to half the requested additive
tolerance; allocating the other half to Krylov convergence changes constants
and logarithms, not the bound. Rayleigh values used in the tests are recomputed
from the returned vectors.
Orthonormalization and Ritz processing
cost \(O(aq_*^2+q_*^3)\) for basis width \(q_*\leq a\), bounded by
\(O(Iq_*)\). This is an eigenvalue/Rayleigh primitive, not an assumed
relative-spectral low-rank oracle for the original product.

Set
\[
 e_* =\theta/8,\qquad \tau=e_*/(512h).
\]
In our parameter range, \(0<\tau\leq1/4096\).
Start with \([\ell,b]=[\ell_0,b_0]\). At its midpoint \(t\), compute
\eqref{eq:ritz} with a fresh random start and form
\(x=J_t^{-\top}y\). If \(q<0\), replace \(\ell\) by \(t\).
If \(q\geq\tau\), replace \(b\) by \(t\). If
\(0\leq q<\tau\), stop and return
\(z=Zx/\sqrt{x^\top Gx}\).
If the bracket first reaches \(b-\ell\leq e_*\ell/4\), make one final
test at \(t=\ell\) and return the vector from that test.

\begin{lemma}[Certified updates and final quality]\label{lem:rayleigh}
On the event that all calls satisfy \eqref{eq:ritz}, the bracket always
contains \(\lambda_1\). The returned vector has unit norm and
\(z^\top Hz\geq(1-e_*)\lambda_1\).
\end{lemma}
\begin{proof}
The inertia of \(K_t\), hence that of \(E_t\), shows that
\(\mu_{\min}(E_t)\geq0\) precisely when \(t\geq\lambda_1\).
Thus \(q<0\) implies \(t<\lambda_1\); and \(q\geq\tau\) implies
\(\mu_{\min}(E_t)\geq q-\tau\geq0\). Both updates are valid.

Whenever \(q\leq\tau\), the witness \(x\) satisfies
\[
 x^\top K_tx\leq\tau x^\top F_tx
   =\tau x^\top K_tx+\tau x^\top\Delta_tx.
\]
Since \(G\succeq I_a/(1+h)\), Lemma~\ref{lem:precond} gives
\begin{equation}\label{eq:witness}
 \frac{x^\top K_tx}{x^\top Gx}\leq64\tau ht,
 \qquad z^\top Hz=t-\frac{x^\top K_tx}{x^\top Gx}
                  \geq t-64\tau ht.
\end{equation}
If the uncertain case occurs, \(q\geq0\) additionally implies
\(\mu_{\min}(E_t)\geq-\tau\). Hence
\(K_t\succeq-\tau F_t=-\tau(K_t+\Delta_t)\), and
\[
 \lambda_1\leq t+
     \frac{42(1+h)\tau ht}{1+\tau}\leq t+45\tau ht.
\]
Combining with \eqref{eq:witness}, \(t\leq3\lambda_1/2\), and
\(\tau=e_*/(512h)\) gives
\(\lambda_1-z^\top Hz\leq(3/2)109 e_*\lambda_1/512<e_*\lambda_1\).

If the final call is at \(t=\ell\), then
\(\mu_{\min}(E_\ell)\leq0\), so \(q\leq\tau\) and
\eqref{eq:witness} applies. Using
\(\lambda_1-\ell\leq e_*\ell/4\) gives
\(\lambda_1-z^\top Hz\leq(e_*/4+e_*/8)\ell<e_*\lambda_1\).
Finally \(\norm{Zx}_2^2=x^\top Gx>0\), since \(G\succ0\) on the
support coordinates. This also proves normalization is valid.
\end{proof}

There are \(O(1+\log(1/e_*))\) tests, since the initial bracket has a
bounded ratio. Every adaptive test uses fresh randomness; condition on its
past and apply the fixed-matrix guarantee with failure probability at most
\(\delta/(4K)\), where a deterministic bound \(K\) on the number of
tests is computed in advance. A union bound costs only logarithms. The total
threshold-test work is \(\tO(I\sqrt{h/\theta}+m^{\omega+\zeta})\).

\section{From a leading Rayleigh vector to spectral low rank}
The following elementary consequence of rank-one interlacing and a Schur
determinant identity avoids any condition number between successive eigenvalues.
It also permits a general-rank extension. The identities themselves are
classical; novelty of this lemma is not asserted.

\begin{lemma}[Relative Nystr\"om deflation]\label{lem:finish}
Let \(H\succeq0\) be \(d\times d\), \(\norm z_2=1\), and
\(s=z^\top Hz>0\). Define
\[
 R=H-\frac{Hz z^\top H}{s}.
\]
Then \(R\succeq0\), \(\rank R=\rank H-1\), and, for
\(1\leq j<d\),
\[
 \lambda_j(R)\leq\frac{\lambda_1(H)}s\lambda_{j+1}(H).
\]
\end{lemma}
\begin{proof}
First suppose \(H\succ0\). In an orthonormal basis starting with \(z\),
write \(H=\left(\begin{smallmatrix}s&b^\top\\b&C\end{smallmatrix}\right)\).
Then \(R=\diag(0,C-bb^\top/s)\). Its positive eigenvalues \(\mu_j\)
satisfy \(\mu_j\geq\lambda_{j+1}(H)\) by rank-one interlacing, and
the Schur determinant identity gives
\[
 \prod_{j=1}^{d-1}\frac{\mu_j}{\lambda_{j+1}(H)}
 =\frac{\det(H)/s}{\prod_{j=1}^{d-1}\lambda_{j+1}(H)}
 =\frac{\lambda_1(H)}s.
\]
Every factor is at least 1, so each is at most their product. This proves
the eigenvalue bound. Replace \(H\) by \(H+\epsilon I\) and let
\(\epsilon\downarrow0\) for the semidefinite case. Finally
\(R=H^{1/2}(I-ww^\top)H^{1/2}\), where
\(w=H^{1/2}z/\sqrt s\) is a unit vector in \(\range H\), proving
positivity and the stated rank drop.
\end{proof}

For the returned \(z\), compute \(s=\norm{M^\top z}_2^2\),
\(v=M^\top z/\sqrt s\), and \(u=Mv\), costing \(O(I)\).
Then \(C=uv^\top=Mvv^\top\) and
\[
 (M-C)(M-C)^\top=H-Hz z^\top H/s.
\]
Since \(s\geq(1-e_*)\lambda_1(H)\), the lemma gives
\(\norm{M-C}_2^2\leq\lambda_2(H)/(1-e_*)
 \leq(1+\theta)^2\lambda_2(H)\).
This also gives exact recovery when \(\lambda_2(H)=0\).
Explicitly form \(C=uv^\top\) in \(O(np)\). This establishes the claimed error
conditional on the good embedding and Krylov events. Their failure
probabilities sum to less than \(\delta\). The embedding application,
Gram formation, support/core work, coarse estimate, tests, final actions,
and output give \eqref{eq:cost}.

\section{General target rank with the same input sketches}
\begin{theorem}[General-rank extension]\label{thm:rank}
For \(1\leq r\leq m\), \(0<\theta\leq h\leq1/16\), there is a
variant returning width at most \(r\), with error at most
\((1+\theta)\sigma_{r+1}(M)\) with probability at least \(1-\delta\),
at total cost
\[
 \tO\!\left(I+G_h+r m^{\omega+\zeta}
       +\frac{r m^2}{\sqrt h}
       +Ir\sqrt{hr/\theta}+\MM(n,r,p)\right).
\]
The algorithm reuses the same two initial embeddings; it does not form a
new pair of tall Gram matrices at each deflation step.
\end{theorem}
\begin{proof}
Maintain orthonormal right vectors \(V_i\), initially empty, and let
\(M_i=M(I-V_iV_i^\top)=AB_i\). At step \(i\), run the threshold
Rayleigh primitive with relative accuracy \(e=\theta/(4r)\) on
\(H_i=M_iM_i^\top\), returning unit \(z_i\).
Here \(e\) is an upper bound on the stage error: calling the original primitive
with parameter \(\theta/r\leq h\) actually gives error
\(\theta/(8r)\leq e\), within its stated parameter range.
Then set
\[
 s_i=z_i^\top H_i z_i,\qquad
 v_i=M_i^\top z_i/\sqrt{s_i},\qquad
 V_{i+1}=[V_i\ v_i].
\]
The new right vector is unit and orthogonal to all previous ones. If the
residual product is zero, detect that from its zero proxy and stop. Otherwise
Lemma~\ref{lem:finish} gives, simultaneously for all \(j\),
\[
 \lambda_j(H_{i+1})\leq(1-e)^{-1}\lambda_{j+1}(H_i).
\]
After \(r\) steps,
\[
 \norm{M-MV_rV_r^\top}_2^2
 \leq(1-e)^{-r}\lambda_{r+1}(H)
 \leq\frac{\lambda_{r+1}(H)}{1-\theta/4}
 \leq(1+\theta)^2\lambda_{r+1}(H),
\]
using Bernoulli's inequality. A zero residual gives exact recovery earlier.

For the arithmetic and adaptive sketch claims, write \(V_i=B^\top C_i\),
\(D_i=BV_i=QC_i\), and \(R_i=I_m-D_iC_i^\top\), so
\(B_i=R_iB=B(I-V_iV_i^\top)\). The initial Loewner inequality for
\(Q\) implies
\[
 (1-h)R_iQR_i^\top\preceq R_i\widetilde QR_i^\top
                         \preceq(1+h)R_iQR_i^\top
\]
for every such adaptive \(R_i\), without another embedding failure event.
At the next step put
\(c=R_i^\top A^\top z_i/\sqrt{s_i}\),
\(v=B^\top c\), and \(d=Bv=Qc\). Orthogonality gives
\(c^\top D_i=0\), and hence
\[
 R_{i+1}=R_i-dc^\top=(I_m-dc^\top)R_i.
\]
The maintained Gram proxy therefore has the rank-two update
\[
 \widetilde Q_{i+1}
 =\widetilde Q_i-dw^\top-wd^\top
                    +(c^\top\widetilde Q_i c)dd^\top,
 \qquad w=\widetilde Q_i c.
\]
This costs \(O(m^2)\), and \(d,v\) cost true input applications.
The fixed support coordinates \(W,L,G\) for \(A\) still apply.
True \(Q_i=Q-D_iD_i^\top\) actions cost \(O(I+mi)=O(I)\).
Every Rayleigh step consequently has the preceding cost with
\(\theta\) replaced by \(\Theta(\theta/r)\). Initial sketch formation
is paid once; small-core factorizations and coarse eigenvalue estimates are
paid \(r\) times. Allocate failure probability \(O(\delta/r)\) to
each stage's Krylov calls. Finally form \(MV_r=AD_r\) in at most
\(O(Ir)\) time and expand the two factors at cost \(\MM(n,r,p)\).
\end{proof}

The initial separate-factor sketches are additional information unavailable
to a matrix-vector oracle for \(M\) alone. The reduced iteration count
therefore does not contradict the matvec lower bounds in the original
direction document.

\section{Sharing the threshold cores across ranks}
\begin{theorem}[One shared core]\label{thm:shared}
The \(r m^{\omega+\zeta}\) term in Theorem~\ref{thm:rank} can be
replaced by \(m^{\omega+\zeta}\). The other terms are unchanged.
\end{theorem}
\begin{proof}
Maintain the explicit \(Q_{0,i}\) by its rank-two update, rather than
reforming \(L\widetilde Q_iL^\top\). If \(a_i=Ld_i\),
\(b_i=Lw_i\), and \(\gamma_i=c_i^\top\widetilde Q_i c_i\), then
\[
 Q_{0,i+1}=Q_{0,i}+
 [a_i\ b_i]\begin{pmatrix}\gamma_i&-1\\-1&0\end{pmatrix}
 [a_i\ b_i]^\top .
\]
Thus \(Q_{0,i}=Q_{0,0}+U_i S_iU_i^\top\), with \(2i\) columns in
\(U_i\) and an invertible block-diagonal \(S_i\), maintained in
\(O(m^2)\) work per stage. The columns need not be independent.

Compute a Frobenius normal form of \(Q_{0,0}\), including its similarity
transformation and inverse, once in \(\widetilde O(m^{\omega+\zeta})\)
field operations \cite{npv}. For any nonpole \(\alpha\), a shifted
companion block solves a new right-hand side by a backward row recurrence
and one division by its characteristic polynomial, in linear work.
Consequently \(B_0^{-1}V\), where \(B_0=\alpha I-Q_{0,0}\), costs
\(O(\MM(m,m,k)+mk)\) for a known \(k\)-column block \(V\).

At a threshold set \(\alpha=(1+32h)t\), \(X=B_0^{-1}U_i\), and
\(J=S_i^{-1}-U_i^\top X\). The actual matrix
\(F_i=B_0-U_iS_iU_i^\top\) is positive definite. The determinant lemma
therefore makes \(J\) invertible whenever \(B_0\) is invertible, and
\[
 F_i^{-1}=B_0^{-1}+XJ^{-1}X^\top.
\]
One application costs \(O(m^2)\). Fix a feasible square exponent
\(9/4<w<\min(5/2,9/4+\zeta)\). Square blocking bounds all preparation,
including forming and inverting \(J\), by
\(O(m^2(1+i)^{w-2})\) per threshold. Its sum over stages is
\[
 \widetilde O(m^2r^{w-1})\leq
 \widetilde O(Ir\sqrt{hr/\theta}).
\]

Run generalized Lanczos on \(F_i^{-1}K_i\) in the \(F_i\) inner
product. It is selfadjoint and analytically equivalent to symmetric
Lanczos on \(F_i^{-1/2}K_iF_i^{-1/2}\), so no new Cholesky factor is
needed. Start with a standard Gaussian \(g\). Since
\(\kappa(F_i)\leq1/h\), its transformed start \(F_i^{1/2}g\)
has normalized squared overlap at least \(h\beta^3/(8a)\) with any fixed unit
minimum eigenvector, except with probability \(\beta\):
use Gaussian anticoncentration for the numerator and
\(\Pr\{\norm g_2^2>2a/\beta\}\leq\beta/2\) for the denominator.
The usual deterministic Chebyshev proof then gives an actual Ritz
quotient within \(\tau/2\) of the minimum eigenvalue in
\(\widetilde O(\tau^{-1/2})\) actions. Only logarithms in \(1/h\)
and \(1/\beta\) are added. All inner products and orthogonalization
costs remain absorbed by the input-action bound.

The auxiliary \(B_0\), unlike \(F_i\), can have poles in later
residual brackets. Test \(a+1\) distinct thresholds in the middle third
of the current bracket. Companion-polynomial evaluation costs \(O(a)\)
per trial, and at most \(a\) trials are poles. Use any nonpole trial;
bracket updates still shrink its width by at least a factor \(2/3\).
This costs \(O(m^2)\) per threshold and introduces no random failure.
For the final query tighten the bracket to
\(b-\ell\leq e_*\ell/4096\), choose a nonpole trial there, and use
Ritz accuracy \(\tau/2\), where \(\tau=e_*/(512h)\).
Even if \(t>\lambda_1\), a top generalized eigenvector gives
\[
 \mu_{\min}\leq
 \frac{\max(0,t-\lambda_1)\norm G_2}{28ht}<\tau/2.
 \]
The actual returned quotient is at most \(\tau\); the witness argument
from Lemma~\ref{lem:rayleigh} gives loss less than
\((1/4096+3/16)e_*\lambda_1<e_*\lambda_1\).
The other sign decisions remain valid with the tighter Ritz tolerance.
Fresh randomness, capped iteration counts, and the previous union bound
complete the claim.
\end{proof}

\section{Adaptive rank and certified absolute error}
The same witnesses permit a stopping rule for a supplied absolute tolerance,
without knowing a suitable rank or any singular value in advance.

\begin{theorem}[Adaptive absolute-error guarantee]\label{thm:absolute}
Given \(\varepsilon>0\), \(0<\theta\leq h\leq1/16\), and
\(0<\delta<1\), an adaptive variant returns factors of a discovered width
\(R\leq m\) and their explicit product \(C\). With probability at least
\(1-\delta\), simultaneously
\[
 \norm{M-C}_2\leq\varepsilon,\qquad
 R\leq r_\varepsilon:=
 \min\{0\leq r\leq m:\ \sigma_{r+1}(M)\leq
                                    \varepsilon/(1+\theta)\}.
\]
For every realized stopping width, its arithmetic cost is
\[
 \tO\!\left(I+G_h+m^{\omega+\zeta}
     +\frac{(R+1)m^2}{\sqrt h}
     +I\sqrt{h/\theta}\,(R+1)^{3/2}
     +np+\MM(n,R,p)\right).
\]
Here \(\MM(n,0,p)=0\); the separate \(np\) term includes writing
the explicit zero output when \(R=0\). No rank-feasibility promise is
needed for the absolute-error guarantee; the spectral counting threshold
controls the returned width and the successful-run cost.
\end{theorem}
\begin{proof}
Reuse the initial embeddings, exact adaptive proxy updates, and shared
core from Theorems~\ref{thm:rank}--\ref{thm:shared}.
At residual stage \(i\geq0\), set
\[
 k_i=\lfloor\log_2(i+1)\rfloor,\qquad
 e_i=\frac{\theta}{32(i+1)(k_i+1)^2}.
\]
Dyadic grouping gives
\[
 \sum_{i\geq0}e_i
 \leq\frac{\theta}{32}\sum_{k\geq0}\frac1{(k+1)^2}
 <\frac{\theta}{16}.
\]
If the current proxy is zero, return the current factors. Otherwise call
the Rayleigh primitive with parameter \(8e_i\leq\theta/4\leq h\).
It returns a unit vector \(z_i\) whose actual quotient
\(s_i=z_i^\top H_i z_i\) satisfies
\[
 (1-e_i)\lambda_1(H_i)\leq s_i\leq\lambda_1(H_i).
\]
Compute \(U_i=s_i/(1-e_i)\). If \(U_i\leq\varepsilon^2\),
return \(C=MV_iV_i^\top\). Otherwise append
\(v_i=M_i^\top z_i/\sqrt{s_i}\), exactly as before. On the good
events the acceptance test certifies the desired residual norm. Every
rejection removes one rank, so the process stops after at most \(m\)
deflations, including exact-zero detection.

Lemma~\ref{lem:finish} and \(s_i\leq\lambda_1(H_i)\) imply
\[
 U_i\leq
 \frac{\lambda_{i+1}(MM^\top)}{\prod_{j=0}^{i}(1-e_j)}
 \leq\frac{\sigma_{i+1}(M)^2}{1-\theta/16}.
\]
The product inequality follows from
\(\prod_j(1-e_j)\geq1-\sum_j e_j\). If
\(i=r_\varepsilon<m\), then
\[
 U_i\leq
 \frac{\varepsilon^2}{(1+\theta)^2(1-\theta/16)}
 \leq\varepsilon^2,
\]
so the algorithm has accepted by that stage. The case
\(r_\varepsilon=m\) follows from rank drop. A zero proxy gives exact
recovery on the common embedding event.
An unexpected primitive failure or zero witness normalization returns the
current factors; this branch is absent on the good events and preserves
the claimed cost in terms of the number of completed deflations.

Give stage \(i\)'s fresh core and threshold Krylov calls total failure
budget at most \(\delta/[4(i+1)(i+2)]\). Conditioning on their past,
the fixed-matrix guarantees and the common embedding event apply;
the summed failure probability is below \(\delta\).
The number of expensive primitive calls is at most \(R+1\), and
\[
 \sum_{i=0}^{R}\sqrt{h/(8e_i)}
 =\widetilde O\!\left(\sqrt{h/\theta}\,(R+1)^{3/2}\right).
\]
All other costs follow from the preceding exact update and output
accounting. The iteration caps also bound runtime on failure outcomes.
\end{proof}

\section{Two levels without forming the fine Gram matrices}
For a distortion parameter \(x\), write \(s_A(x),s_B(x)\) for the
embedding dimensions from Section 1 and define
\[
 \mathcal J_x=m\bigl(s_A(x)+s_B(x)\bigr)+m^2.
\]
This is the cost of applying the two sketched Gram operators and the
small coordinate transformations.

\begin{theorem}[Matrix-free fine proxies]\label{thm:twolevel}
For \(1\leq r\leq m\), \(0<\theta\leq h\leq H/4\), and
\(H\leq1/32\), there is an algorithm returning width at most \(r\)
and the explicit product \(C\), such that
\[
 \Pr\{\norm{M-C}_2\leq(1+\theta)\sigma_{r+1}(M)\}\geq1-\delta.
\]
Its total arithmetic cost is
\begin{equation}\label{eq:twolevel}
 \tO\!\left(I+G_H+m^{\omega+\zeta}
    +\frac{rm^2}{\sqrt H}
    +Ir\sqrt{hr/\theta}
    +\mathcal J_h r\sqrt{Hr/\theta}
    +\MM(n,r,p)\right).
\end{equation}
Only the coarse \(H\)-accuracy sketches are multiplied into dense Grams.
The finer \(h\)-accuracy sketches remain in factored form.
\end{theorem}
\begin{proof}
First consider rank one. Form coarse \(P_H,Q_H\) and use the preceding
support construction at distortion \(H\). It supplies the coordinate maps,
the initial bracket \([\ell_0,b_0]\), and
\begin{gather*}
 F_H(t)=(1+32H)tI-Q_0,\\
 22HtI\preceq F_H(t)-K_t\preceq42HtI.
\end{gather*}
The coarse norm estimate costs \(\widetilde O(m^2/\sqrt H)\).
Compute finer sketch factors
\[
 A_h=S_hA,\quad B_h=BT_h^\top,\quad
 P_h=A_h^\top A_h,\quad Q_h=B_hB_h^\top,
\]
where the last two matrices are used only as operators. In these same
coarse coordinates define
\[
 \widehat G=W^\top P_hW,\quad \widehat Q=LQ_hL^\top,\quad
 \widehat K_t=t\widehat G-\widehat G\widehat Q\widehat G,\quad
 F_h(t)=\widehat K_t+32htI .
\]
Every such action costs \(O(\mathcal J_h)\); neither a fine dense Gram
nor a tall transformed factor is formed. The fine Loewner event implies
\(\norm{\widehat G-G}_2\leq h\norm G_2\) and
\(\norm{\widehat Q-Q'}_2\leq h\norm{Q'}_2\).
Using \(\norm G_2\leq16/15\) and \(\norm{Q'}_2\leq2t\) gives
\[
 \norm{\widehat K_t-K_t}_2\leq10ht,\qquad
 22htI\preceq F_h(t)-K_t\preceq42htI.                 \tag{*}
\]
Keeping the mass matrix \(\widehat G\) in this formula is essential.

\paragraph{A fine bracket using only fine-factor actions.}
Let \(\lambda=\norm M_2^2\) and
\(\lambda_h=\norm{A_hB_h}_2^2\). The fine embeddings give
\((1-h)^2\lambda\leq\lambda_h\leq(1+h)^2\lambda\).
The supplied \(P_H,Q_H\) are also \(2H\)-relative Gram proxies for
the fine factors: their Loewner ratios lie between
\((1-H)/(1+h)\) and \((1+H)/(1-h)\), within \([1-2H,1+2H]\).
Apply the already proved Rayleigh primitive to \(A_h,B_h\), using
these supplied cores, distortion \(2H\), and requested parameter \(h\).
It needs no additional embedding or fine Gram. Its true-factor actions
cost \(\mathcal J_h\), and it returns an actual quotient \(s_h\) with
\((1-h/8)\lambda_h\leq s_h\leq\lambda_h\), at cost
\[
 \tO\bigl(\mathcal J_h\sqrt{H/h}+m^{\omega+\zeta}
                                  +m^2/\sqrt H\bigr).
\]
Intersect \([\ell_0,b_0]\) with
\[
 \left[\frac{s_h}{(1+h)^2},
       \frac{s_h}{(1-h/8)(1-h)^2}\right].
\]
On the good events this is a nonempty bracket for \(\lambda\), and its
lower endpoint is at least \((1-33h/8)\lambda\).
Every threshold in it satisfies \(K_t\succeq-5htI\), hence
\(F_h(t)\succeq16htI\). Also, since \(h\leq H/4\), (*) gives
\[
 0\preceq F_H(t)-F_h(t)\preceq42HtI,\qquad
 F_H(t)\preceq\bigl(1+42H/(16h)\bigr)F_h(t).
\]

\paragraph{A fixed polynomial inverse root.}
For the analysis, factor \(F_H(t)=JJ^\top\) and put
\[
 D=J^{-1}F_h(t)J^{-\top},\qquad
 a=\frac{h}{3H},\qquad \eta=\frac h{1000}.
\]
The preceding comparison gives \(aI\preceq D\preceq I\).
A polynomial \(p\) with
\(\left|\sqrt{x}\,p(x)-1\right|\leq\eta\) on \([a,1]\)
can be applied in \(\widetilde O(a^{-1/2})\) actions of \(D\).
For completeness an explicit arithmetic construction is
\[
 \rho=\frac{1-\sqrt a}{1+\sqrt a},\quad
 c=\frac{(1+\sqrt a)^2}{4},\quad
 T=\frac{(1+a)I-2D}{1-a},\quad
 p(D)=c^{-1/2}\sum_{j=0}^{k}\rho^j P_j(T),
\]
where \(P_j\) are the Legendre polynomials. Their generating function
\cite{legendre} and \(\norm{P_j(T)}_2\leq1\) bound the tail by
\(\rho^{k+1}/\sqrt a\). Choose
\(\rho^{k+1}\leq\eta\sqrt a\), so
\(k=O(a^{-1/2}\log(1/(\eta\sqrt a)))\).
The three-term Legendre recurrence uses \(k\) operator actions,
\(O(k)\) scalar operations, and \(O(km)\) vector work. Repeated
multiplication by \(\rho\) finds \(k\); no logarithm, trigonometric,
or coefficient-integration oracle is needed.

Define analytically
\[
 R=J^{-\top}p(D),\quad E_t=R^\top K_tR,\quad
 \overline F=R^{-\top}R^{-1}=Jp(D)^{-2}J^\top .
\]
Neither \(\overline F\) nor \(p(D)^{-1}\) is formed. Relative polynomial
error implies
\(-4\eta F_h\preceq\overline F-F_h\preceq4\eta F_h\).
Since \(F_h\preceq F_H\preceq2tI\), it follows that
\[
 21htI\preceq\overline F-K_t\preceq43htI,\qquad
 \overline F\succeq15htI.
\]
Thus \(E_t\preceq I\); using \(K_t\succeq-5htI\) also gives
\(E_t\succeq-I/3\). The old actual-witness threshold argument applies:
for \(x=Ry\), a quotient \(q\leq\tau\) gives
\(x^\top K_tx/(x^\top Gx)\leq64\tau ht\); an uncertain sign gives
\(\lambda\leq t+45\tau ht\). The constants hold since \(H\leq1/32\).
With \(\tau=e_*/(512h)\), \(e_*=\theta/8\), the returned original
Rayleigh vector is \((1-e_*)\)-accurate. A single \(E_t\) action costs
\(\widetilde O(I+\mathcal J_h\sqrt{H/h})\).
The total action cost is consequently
\[
 \tO\bigl(I\sqrt{h/\theta}
                  +\mathcal J_h\sqrt{H/\theta}\bigr).
\]
This absorbs the fine-bracket cost because \(\theta\leq h\).
The Nystr\"om lemma supplies the spectral approximation and explicit output.

\paragraph{Sharing and deflation.}
There is no need to compute \(J\) in the algorithm. Use the shared exact
\(F_H^{-1}\) from Theorem~\ref{thm:shared}, and let
\(\mathcal T=F_H^{-1}F_h\), selfadjoint in the \(F_H\) inner product.
Run generalized Lanczos on
\[
 p(\mathcal T)F_H^{-1}K_t p(\mathcal T)
\]
in that inner product. Its symmetric coordinates are exactly \(E_t\).
For its \(F_H\)-unit witness \(x\), put \(w=p(\mathcal T)x\);
its actual quotient is \(w^\top K_tw\), and the final original
vector is \(Zw/\sqrt{w^\top Gw}\).
Gaussian starts incur only the logarithmic penalty
\(\kappa(F_H)=O(1/H)\) proved above.
The nonpole query rule and tighter final bracket from the shared-core
proof still work, now with \(\overline F\succeq15htI\).
In particular their terminal minimum eigenvalue is below \(\tau/2\).

For general rank, use stage error \(\theta/(8r)\) and the same
Nystr\"om deflation. Both pairs of original embedding events survive
every adaptive congruence \(Q\mapsto R_iQR_i^\top\).
The coarse cores have the previous rank-two updates. Fine residual
actions apply \(R_iQ_hR_i^\top\) through the stored fine factor and
the rank-\(i\) transformations, still in \(O(\mathcal J_h)\) work.
The same shared Frobenius form supports both the \(H\) and \(2H\)
coarse-proxy threshold shifts. Its per-stage preparation is absorbed
as in Theorem~\ref{thm:shared}. Summing the action and coarse-norm
costs gives \eqref{eq:twolevel}. Fresh Krylov seeds conditional on the
past, and a union bound over the four initial embeddings and all
stage calls, give success probability \(1-\delta\).
\end{proof}

\begin{corollary}[More accurate quadratic output]\label{cor:twolevel}
Let \(n=p=N\), \(m=\lceil N^{50/59}\rceil\),
\[
 \theta=m^{-11/25},\qquad H=m^{-7/100},\qquad
 h=m^{-17/200},\qquad\delta=1/100.
\]
Take \(N\) sufficiently large for the parameter inequalities.
For every fixed target rank \(r\) and every fixed real factor pair,
the two-level algorithm returns the explicitly expanded \(C\) in total
\(O(N^2)\) arithmetic operations, with error at most
\((1+\theta)\sigma_{r+1}(M)\) with probability at least \(99/100\).
Construction takes \(o(N^2)\).
\end{corollary}
\begin{proof}
In powers of \(m\), the input has exponent \(109/50\), the output
\(59/25\), the coarse sketch row count \(57/50\), and the fine row
count \(117/100<59/50\). The three potentially leading construction
terms have exponents
\[
 G_H:\ \frac{3362}{1425},\qquad
 I\sqrt{h/\theta}:\ \frac{943}{400},\qquad
 \mathcal J_h\sqrt{H/\theta}:\ \frac{471}{200}.
\]
Their margins below \(59/25\) are respectively \(1/1425\),
\(1/400\), and \(1/200\).
The core exponent \(9/4\) and coarse norm exponent \(407/200\)
are also strictly smaller. Fix sufficiently small exponent slack;
all logarithms are absorbed in these strict margins.
For fixed \(r\), conventional output expansion costs \(O(N^2r)=O(N^2)\).
\end{proof}

At this accuracy the named exact-Gram bound is still
\(N^{7056/3481+o(1)}\), and ordinary Krylov gives
\(N^{120/59+o(1)}\). This also improves the optimized one-level
tradeoff: a one-level distortion \(m^{-a}\) would need \(a\geq2/25\)
for its true-action term to fit \(N^2\), but its Gram exponent would
then be at least
\[
 \frac{54}{25}+\frac{25}{116}>
 \frac{59}{25}.
\]
Larger sketch row counts only increase that setup bound; clipping at
the original dimension restores the exact-Gram comparison.

The absolute-error stopping rule from Theorem~\ref{thm:absolute}
applies to the two-level primitive as well. Its cost replaces the
two rank-dependent action terms in \eqref{eq:twolevel} by
\[
 \bigl(I\sqrt{h/\theta}+\mathcal J_h\sqrt{H/\theta}\bigr)(R+1)^{3/2},
\]
and \(r\) in the coarse-norm term by \(R+1\), with an additional
\(np\) write term when \(R=0\). The width and error guarantees are
unchanged. In the displayed regime, constant \(r_\varepsilon\) gives,
on the same event of probability at least \(1-\delta\), adaptive
absolute error and \(O(N^2)\) total time, without supplying that rank
to the algorithm.

\section{A fixed number of finer levels}
The same polynomial construction can be nested. Uniform inflation constants
and explicit validity windows control both the inverse applications and the
intermediate norm brackets.

\begin{theorem}[Fixed-depth tradeoff]\label{thm:multilevel}
Fix an integer \(L\geq1\) independently of the dimensions and accuracy.
Let \(h_0\leq1/1024\), \(h_j\leq h_{j-1}/8\) for \(1\leq j\leq L\),
and \(0<\theta\leq h_L\). A rank-one approximation with the same
\((1+\theta)\sigma_2(M)\) guarantee has total cost
\[
 \tO_L\!\left(I+G_{h_0}+m^{\omega+\zeta}+\frac{m^2}{\sqrt{h_0}}
       +I\sqrt{h_L/\theta}
       +\sum_{j=1}^L\mathcal J_{h_j}\sqrt{h_{j-1}/\theta}
       +np\right).
\]
Constants and the exponent of the suppressed logarithms may depend on
fixed \(L\). The statement makes no claim for input-dependent growing depth.
\end{theorem}
\begin{proof}
Use the \(h_0\)-Gram support coordinates. Let \(K\) be the true
threshold pencil, let \(K_j\) be its proxy using the level-\(j\)
factored sketches, and set \(K_0=tI-Q_0\). Define
\(F_j=K_j+128h_jtI\).
Call a threshold level-\(j\) valid when
\(\lvert t-\lambda\rvert\leq16h_j\lambda\), where
\(\lambda=\norm M_2^2\). This condition is used in the analysis; the
algorithm uses the brackets below.
On such a window, \(\norm{Q'}_2\leq2t\), so the earlier proxy
expansion still gives \(\norm{K_j-K}_2\leq10h_jt\). In particular
\[
 K\succeq-17h_jtI,\quad
 118h_jtI\preceq F_j-K\preceq138h_jtI,\quad
 F_j\succeq100h_jtI,\quad \norm{F_j}_2\leq4t.
\]
These estimates hold on the full validity window, rather than just
the original narrower bracket.
The error estimate for a finer \(K_j\) also holds in every coarser
validity window: its derivation only needs \(\norm{Q'}_2\leq2t\).

Factor \(F_0=J_0J_0^\top\), set \(R_0=J_0^{-\top}\), and
inductively maintain the analytical effective preconditioner
\(\overline F_{j-1}=R_{j-1}^{-\top}R_{j-1}^{-1}\) with
\[
 117h_{j-1}tI\preceq\overline F_{j-1}-K
                      \preceq139h_{j-1}tI,\qquad
 \overline F_{j-1}\succeq99h_{j-1}tI .
\]
At a level-\(j\) valid threshold, the separation of consecutive
distortions gives
\[
 0\preceq\overline F_{j-1}-F_j\preceq139h_{j-1}tI,\qquad
 \frac{h_j}{2h_{j-1}}I\preceq
       D_j:=R_{j-1}^\top F_jR_{j-1}\preceq I.
\]
Use the explicit Legendre polynomial with relative inverse-root error
\(\eta_j=h_j/1000\) on this interval and set
\(R_j=R_{j-1}p_j(D_j)\).
Its effective preconditioner differs from \(F_j\) in norm by at most
\(4\eta_j\norm{F_j}_2\leq16h_jt/1000\).
Thus the displayed induction bounds hold at level \(j\), and
\(\operatorname{spec}(R_j^\top K R_j)\subset[-1,1]\).
The approximation errors are reset relative to \(F_j\) at every
level; no unproved sum of nested solve errors is used.

It remains to construct all required validity windows. The original
coarse norm bracket lies within
\([(1-6h_0)\lambda,(1+6h_0)\lambda]\).
Suppose such a true-\(\lambda\) bracket is available at level \(j-1\).
For \(\lambda_j=\norm{A_jB_j}_2^2\), expand its endpoints by
\((1-h_j)^2\) and \((1+h_j)^2\), respectively.
Every resulting threshold is within \(8h_{j-1}\lambda\) of
\(\lambda\), hence is valid at level \(j-1\).
Use \(R_{j-1}^\top K_jR_{j-1}\) to estimate \(\lambda_j\).
Its effective-preconditioner difference is between
\(115h_{j-1}tI\) and \(141h_{j-1}tI\), so its spectrum is in
\([-1,1]\). The true mass in this step is the fine
\(G_j=W^\top P_jW\), whose eigenvalues lie between \(0.99\) and \(1.01\).
The threshold witness and uncertainty inequalities both hold with
constant \(256\). Use tolerance
\(\tau=e/(2048h_{j-1})\) and \(e=h_j/8\).
The uncertain-case loss is at most
\((3/2)(256+256)e/2048<e\).
The final lower-endpoint query obeys the same witness bound.
This produces an actual quotient \(s_j\) satisfying
\((1-h_j/8)\lambda_j\leq s_j\leq\lambda_j\).
Translate it to the true-\(\lambda\) bracket
\[
 \left[\frac{s_j}{(1+h_j)^2},
       \frac{s_j}{(1-h_j/8)(1-h_j)^2}\right].
\]
It lies within \([(1-5h_j)\lambda,(1+5h_j)\lambda]\), establishing
the induction. These refinements use only fine-product actions.
Finally use \(R_L^\top K R_L\) on the last bracket, with
\(e=\theta/8\), \(\tau=e/(2048h_L)\), and the same \(256\) bounds.
Nystr\"om finishing gives the required spectral error.

If \(c_j\) is the cost of applying \(R_j\) or its transpose at a
fixed threshold, the polynomial degrees give
\[
 c_0=O(m^2),\qquad
 c_j=\tO_L\bigl((\mathcal J_{h_j}+c_{j-1})
                                     \sqrt{h_{j-1}/h_j}\bigr).
\]
Expanding this recurrence and multiplying by
\(\widetilde O(\sqrt{h_L/\theta})\), the final true-operator count,
gives the stated sum. Each intermediate norm-refinement cost is
absorbed because \(\theta\leq h_L\leq h_j\).
The propagated \(m^2\) cost is absorbed by the first sketch term.
Only logarithmically many coarse factorizations are needed per level.
All levels use fixed polynomials at each threshold; fresh random
starts and a union bound over the fixed embedding pairs and adaptive
calls give the stated confidence. The usual capped failure exits
handle zero products and invalid normalizations.
\end{proof}

\begin{corollary}[Two finer sketch levels]\label{cor:multilevel}
Let \(n=p=N\), \(m=\lceil N^{50/59}\rceil\), \(\delta=1/100\), and
\[
 \theta=m^{-56/125},\qquad
 h_0=m^{-7/100},\quad h_1=m^{-17/200},\quad h_2=m^{-89/1000}.
\]
Take \(N\) sufficiently large for the parameter inequalities.
For every fixed rank \(r\), every fixed factor pair admits the stated
randomized \((1+\theta)\sigma_{r+1}(M)\) guarantee in total
\(O(N^2)\) arithmetic operations, including explicit output.
The construction cost is \(N^{2-1/2360+o(1)}\).
\end{corollary}
\begin{proof}
For rank one the exponents in powers of \(m\) are
\[
 G_{h_0}:\frac{3362}{1425},\quad
 \mathcal J_{h_1}\sqrt{h_0/\theta}:\frac{2359}{1000},\quad
 \mathcal J_{h_2}\sqrt{h_1/\theta}:\frac{4719}{2000},\quad
 I\sqrt{h_2/\theta}:\frac{4719}{2000}.
\]
Their maximum is \(4719/2000=59/25-1/2000\); converting to
powers of \(N=m^{59/50+o(1)}\) gives \(2-1/2360\).
The core and coarse-norm terms are smaller.
Fix multiplication slack below \(23/114000\) in \(m\)-based exponents,
the gap between \(4719/2000\) and \(3362/1425\).
Output expansion is \(O(N^2)\).
For any fixed \(r\), apply the proved Nystr\"om deflation with stage
accuracy \(O(\theta/r)\), reusing every embedding by exact congruences.
Stage counts, the finer precision, and coarse factorizations only change
constants depending on \(r\). The exponent and the guarantee follow.
\end{proof}

The adaptive absolute-error rule also applies to this primitive.
When \(r_\varepsilon\) is constant, with probability at least \(1-\delta\)
both the absolute-error guarantee and the quadratic total cost hold,
without a supplied rank. The strongest
checked general-product Gram comparison remains
\(N^{7056/3481+o(1)}\); ordinary Krylov here gives
\(N^{601/295+o(1)}\). The theorem concerns real arithmetic and
asymptotic cost; it does not establish a finite-size numerical speedup.

\section{The one-level quadratic-output instance}
\begin{corollary}[Explicit output in quadratic time]\label{cor:output}
Let \(n=p=N\), \(m=\lceil N^{50/59}\rceil\),
\(\theta=m^{-25/59}\), \(h=m^{-7/100}\), and \(\delta=1/100\).
For every fixed pair of real factors of these dimensions, the algorithm
returns \(C=uv^\top\), including every entry of \(C\), in total
\(O(N^2)\) arithmetic operations with
\[
 \Pr\{\norm{M-C}_2\leq(1+\theta)\sigma_2(M)\}\geq99/100.
\]
Its construction takes \(o(N^2)\) time.
\end{corollary}
\begin{proof}
In powers of \(m\), the outer dimension is \(m^{59/50+o(1)}\),
the input has exponent \(109/50\), and the output has exponent
\(59/25\). The sketch has \(\tO(m^{57/50})\) rows, strictly fewer
than the outer dimension by a polynomial margin. Reference \cite{amanov}
therefore gives
\[
 G_h\leq m^{3362/1425+o(1)},\qquad
 I\sqrt{h/\theta}=m^{27811/11800+o(1)}.
\]
The margins below the output exponent are, respectively,
\[
 \frac{59}{25}-\frac{3362}{1425}=\frac1{1425},\qquad
 \frac{59}{25}-\frac{27811}{11800}=\frac{37}{11800}.
\]
The square core has exponent \(9/4\), and the coarse eigenvalue estimate
has exponent \(2+7/200=407/200\), also strictly below \(59/25\).
Choose fixed multiplication exponent slack below \(1/1425\), leaving
room to absorb all logarithms. Every construction term is \(o(N^2)\).
The final rank-one expansion costs \(O(N^2)\).
\end{proof}

For every fixed target rank \(r\), Theorem~\ref{thm:rank} gives the same
\(O(N^2)\) total-time consequence, with \(\sigma_2\) replaced by
\(\sigma_{r+1}\). The factors \(r\) and \(\sqrt r\) do not change
the exponent, and conventional expansion costs \(O(N^2r)=O(N^2)\).

In \(N\)-based exponents, the two leading construction bounds are
\[
 N^{6724/3363+o(1)}\quad\text{and}\quad
 N^{27811/13924+o(1)},
\]
both below 2. For the same dimensions and accuracy, the named exact-Gram
and ordinary-Krylov comparison bounds both have exponent
\[
 \frac{7056}{3481}\approx2.027004.
\]
The named exact-product exponent is \(25281/11800\approx2.142458\).
Thus the output-optimal bound saves \(N^{94/3481-o(1)}\) over those
general-product Gram/Krylov bounds. It changes the parameter regime rather
than proving the original \(m=N^{8/9},\theta=1/m\) target.

For fixed absolute tolerance \(\varepsilon>0\), the same conclusion gives
\(\norm{M-C}_2\leq\varepsilon\) with probability at least \(99/100\)
whenever \((1+\theta)\sigma_2(M)\leq\varepsilon\). No factor-energy,
conditioning, or spectral-gap promise is imposed. Without this rank-one
feasibility promise, the relative guarantee still holds on every input;
it does not assert that an arbitrary product admits rank-one absolute error
\(\varepsilon\).

\section{A comparison at the original inner dimension}
Take \(n=p=N\), \(m=\lceil N^{8/9}\rceil\), but now choose
\[
 \theta=m^{-1/2},\qquad h=m^{-3/50},\qquad\delta=1/100.
\]
This differs from the direction document's harder \(\theta=1/m\) example.
The rectangular bound in \cite{amanov} gives, in powers of \(m\),
\[
 I=m^{17/8+o(1)},\qquad s_A,s_B=m^{28/25+o(1)},\qquad
 G_h\leq m^{6561/2800+o(1)},
\]
where \(6561/2800\approx2.343214\). Meanwhile
\[
 I\sqrt{h/\theta}=m^{469/200+o(1)}=m^{2.345+o(1)},\quad
 np=m^{9/4+o(1)},\quad m^2/\sqrt h=m^{203/100+o(1)}.
\]
Thus the candidate total bound is
\[
 N^{469/225+o(1)}\quad(469/225\approx2.084444).
\]
The exact-Gram comparison remains
\(N^{169/81+o(1)}\approx N^{2.086420}\); ordinary Krylov with this
accuracy gives \(N^{19/9+o(1)}\), and exact rectangular multiplication
gives \(N^{625/288+o(1)}\). The saving over the named Gram route is
\(N^{4/2025-o(1)}\), a small positive exponent. Choose the fixed exponent
slack sufficiently small to preserve it.

The relative guarantee in this second example implies absolute error when
\((1+\theta)\sigma_2(M)\leq\varepsilon\), exactly as in the original
rank-one comparison. No factor-energy, spectral-gap, or conditioning promise
has been added. These are comparisons against the audited general-product
routes. Specialized covariance algorithms can be faster on
\(B=A^\top\); for example, accelerated stochastic PCA \cite{garber}
exploits explicit covariance samples unavailable at that cost for arbitrary
two-factor products. We do not claim to dominate every structured subclass.
A complete state-of-the-art novelty claim requires a broad prior-work review.
The checks completed so far are summarized below. The theorem has not been
machine formalized or subjected to an external referee review.

\section{Scope of the prior-work comparison}
The claimed new contribution is the analyzed product-specific tradeoff, not
the invention of preconditioned Lanczos, embeddings, or Nystr\"om extraction.
The following primary-source checks found no earlier theorem yielding the
stated total bound on arbitrary two-factor inputs.

\paragraph{Related eigenvalue methods.}
Classical shifted-pencil Lanczos \cite{templates} supplies the mechanism.
The general gap-free bounds of \cite{genev} retain the full inverse-square-root
accuracy factor. EPSI \cite{epsi} is closer in spirit, but the checked
convergence guarantees are local and gap dependent. RLOBPCG \cite{rlobpcg}
treats smallest singular vectors of a supplied tall matrix under a gap
condition; it does not supply a cheap factor for the indefinite threshold
matrix \(tG-GQ'G\). The current survey \cite{survey} and its closest
references were checked in addition to these original papers.

\paragraph{Oversampling and fast rectangular multiplication.}
Large-block finite-Schatten algorithms \cite{kw} and direct spectral
oversampling bounds \cite{tropp}, including blocks wider than the target
rank, require particular care. Our audit of their published guarantees
found no uniform spectral improvement that subsumes the bound here: the
certified pass count retains the inverse-square-root accuracy dependence,
or a finite-Schatten conversion requires a full-width block. This is a
comparison of proved guarantees, not a lower bound against every possible
oversampled algorithm or a claim that every input realizes their worst case.

\paragraph{Ridge and stochastic methods.}
The multilevel methods in \cite{msp} include algorithms for supplied
rectangular features that avoid forming their Gram matrices. Their
positive-ridge systems nevertheless require access to those positive
Gram features. The signed, noncommuting threshold pencil here has no
supplied cheap positive Gram root; substituting such a root must include
its construction cost. The checked theorems do not supply our threshold
oracle or the resulting two-factor spectral bound. Recursive
preconditioning itself is established prior art.
Stochastic covariance methods \cite{garber} remain important stronger
specializations when actual covariance samples are cheap. Extending those
particular sample-cost bounds directly to general \(AB\) would omit
feature formation or true-Gram work.

Product-entry sampling, COD, stable-rank AMM, structured Gram recovery,
and recent output-Frobenius algorithms were also audited in the supporting
research records. Their stated promises or error terms do not supply the
uniform relative spectral guarantee and total bound here. These non-findings
support a candidate improvement over the located general-product literature;
they are not a formal certificate of exhaustive novelty.

\begin{thebibliography}{99}
\bibitem{mm} C. Musco and C. Musco. \emph{Randomized Block Krylov Methods
for Stronger and Faster Approximate Singular Value Decomposition} (2015).
\url{https://arxiv.org/abs/1504.05477}.
\bibitem{cnw} M. Cohen, J. Nelson, D. Woodruff. \emph{Optimal Approximate
Matrix Product in Terms of Stable Rank} (2016), Appendix A.2.
\url{https://arxiv.org/abs/1507.02268}.
\bibitem{amanov} A. Amanov, M. van den Berg, J. Zuiddam.
\emph{Multiplying Matrices via Polynomials: Exposition and Barriers}
(2026 preprint), Theorem 1.2.
\url{https://arxiv.org/abs/2610.30054}.
\bibitem{epsi} R. Xu, Y. Lu. \emph{What Is a Sketch-and-Precondition
Derivation for Low-Rank Approximation? Inverse Power Error or Inverse Power
Estimation?} (2025 preprint).
\url{https://arxiv.org/abs/2502.07993}.
\bibitem{jyz} R. Jiang, M.-C. Yue, Z. Zhou.
\emph{An Accelerated First-Order Method with Complexity Analysis for Solving
Cubic Regularization Subproblems}. Computational Optimization and Applications
79 (2021), Lemma 1.
\url{https://rjjiang.github.io/papers/COAP2021.pdf}.
\bibitem{templates} A. Knyazev. \emph{Preconditioned Lanczos Methods}, in
\emph{Templates for the Solution of Algebraic Eigenvalue Problems} (2000).
\url{https://netlib.org/utk/people/JackDongarra/etemplates/node415.html}.
\bibitem{garber} D. Garber et al. \emph{Faster Eigenvector Computation via
Shift-and-Invert Preconditioning} (2016).
\url{https://arxiv.org/abs/1605.08754}.
\bibitem{genev} Z. Allen-Zhu, Y. Li. \emph{Doubly Accelerated Methods for
Faster CCA and Generalized Eigendecomposition} (2017).
\url{https://arxiv.org/abs/1607.06017}.
\bibitem{rlobpcg} E. Epperly, T. Park, Y. Nakatsukasa.
\emph{Fast, High-Accuracy, Randomized Nullspace Computations for Tall Matrices}
(2026 preprint).
\url{https://arxiv.org/abs/2602.16797}.
\bibitem{survey} E. Epperly, D. Kressner, Y. Nakatsukasa.
\emph{Randomized Algorithms in Numerical Linear Algebra for Scientific Computing}
(2026 preprint).
\url{https://arxiv.org/abs/2610.38157}.
\bibitem{kw} P. Kacham, D. Woodruff. \emph{Faster Algorithms for Schatten-p
Low Rank Approximation} (2024).
\url{https://arxiv.org/abs/2407.11959}.
\bibitem{tropp} J. Tropp. \emph{Analysis of Randomized Block Krylov Methods},
Theorem 2.5.
\url{https://tropp.caltech.edu/reports/Tro18-Analysis-Randomized-TR.pdf}.
\bibitem{msp} M. Derezi\'nski, C. Musco, J. Yang. \emph{Faster Linear
Systems and Matrix Norm Approximation via Multi-level Sketched Preconditioning}
(2025).
\url{https://arxiv.org/abs/2405.05865}.
\bibitem{npv} V. Neiger, C. Pernet, G. Villard.
\emph{Computing Krylov Iterates in the Time of Matrix Multiplication}
(2024), Corollary 5.2.
\url{https://arxiv.org/abs/2402.07345}.
\bibitem{legendre} NIST Digital Library of Mathematical Functions.
Legendre generating function, Equation 18.12.11.
\url{https://dlmf.nist.gov/18.12.E11}.
\end{thebibliography}
\end{document}
